
\section{Customise the data reduction}
\label{sec:edps_parameters}
Each step of the reduction cascade is organised in the so-called "tasks". In other words, a task represent a step in the data reduction. Each task is defined by input files and the recipe to execute. It is possible to define the reduction cascade up to a certain step, by selecting the "target task", or a "meta-target" (i.e. a defined group of tasks), so that the data reduction stops at a certain point.

In order to display the tasks for a certain workflow type:
\begin{verbatim}
  edps -w espresso.espresso_wkf -lt
 \end{verbatim}  
The output list shows the tasks grouped by meta-target.
While a task can run only one recipe, a recipe can be associated to many tasks.

In order to execute the reduction up to a certain step, one has to indicate the "target task", or a meta-target (this will select multiple "target tasks", i.e. those that are within the same meta-target group). Type:
\begin{verbatim}
  edps -w espresso.espresso_wkf -t <TASK>
 \end{verbatim}  
for a single "target task", or:
 \begin{verbatim}
  edps -w espresso.espresso_wkf -m <meta-target>
 \end{verbatim}  
for multiple target tasks, within meta-target group.
 More information on Section \ref{sec:edps_cascade}.


\subsection{Recipe parameters}
\label{sec:edps_recipe_parameters}

To process data with specific value for a parameter, which is executed by a recipe of a given task, type on a terminal where the edps Python environment is active (single line command):
\begin{verbatim}
  edps -w espresso.espresso_wkf -i <input_directories> 
       -rp <TASK> <PARAMETER> <VALUE> 
       -o <output_directory>
 \end{verbatim}  
 
Where TASK is the task name that runs the recipe we want to change the parameter for; PARAMETER is the parameter name, and VALUE is the value we want to use. The PARAMETER name must give the full name, which includes the instrument name and the recipe name (see Section \ref{sec:edps_default_params} on how to display the available recipe parameters). 

For example, the command (single line):

\begin{verbatim}
  edps -w  espresso.espresso_wkf -i <input_directories> 
       -rp object espdr.espdr_sci_red.cosmic_detection_sw 1
       -o <output_directory>
 \end{verbatim}
 
instructs the task {\tt object} to activate the Laplacian Cosmic Ray detection algorithm when running the recipe {\tt espdr\_sci\_red}.

To change more than one parameters, just add other {\tt -rp} lines to the example above. If many parameters have to be configured, it might be convenient to load them from a configuration file (see Section \ref{sec:edps_parameters_file})

For a given instrument workflow, the full list of options is described in the corresponding pipeline manuals and data reduction tutorials.

{\it Note:} The recipe parameters defined through the command line as above override:
\begin{itemize}
    \item default recipe values;
    \item configuration files with specified values;
    \item values that are hard-coded or configured automatically by the workflow (e.g., that depend on input data).
\end{itemize}


\subsection{Workflow parameters: select the most appropriate data reduction strategy}
\label{sec:edps_workflow_parameters}

Some instrument workflows can reduce the data in different ways, depending on the science needs. Certain steps of the reduction can be avoided or executed instead of others, or some calibrations can be ignored despite being present in the input data directory. 

Some of the data reduction strategies are hard-coded in the workflow, in the sense that the data reduction cascade depends on the properties of the data we want to process. On the other hand, some other strategies can be controlled by so called "workflow parameters." In other words, workflow parameters are not directly associated to the recipes but they define the strategy of the data reduction and the reduction chain.
In the case of ESPRESSO, there is only one reduction strategy and no workflows parameters are available. In the case of the KMOS workflow, the following workflow parameters are available:
\begin{itemize}
    \item molecfit: 'standard' (default), 'science', 'false'.
    \item use\_sky\_flats: 'false' (default), 'true'.
    \item qc0: 'false' (default), 'true'.
\end{itemize}

For example, to use the sky flats for illumination correction instead of using the lamp flats as default strategy, type (single line command):
\begin{verbatim}
  edps -w kmos.kmos_wkf -i <input_directories> 
       -wp use_sky_flats 'true'
       -o <output_directory> 
 \end{verbatim}

For a given instrument workflow, the full list of options is described in the corresponding pipeline manuals and data reduction tutorials.


To specify more workflow parameters, add as many {\tt -wp} lines as needed. If a large number of workflow parameters have to be specified, it might be convenient to  load them from a configuration file (see Section \ref{sec:edps_parameters_file})

{\it Note:} Workflow parameters defined  through the command line as above, override:
\begin{itemize}
    \item values defined in configuration files;
    \item values that are hard-coded or configured automatically by the workflow (e.g., that depends on input data).
\end{itemize}



\subsection{Display default parameter values}
\label{sec:edps_default_params}
To display what are the recipe parameters full names and their values for the task {\tt object} used in the default parameter set (note: the ESPRESSO pipeline must be installed):

\begin{verbatim}
 edps -w espresso.espresso_wkf -p object
\end{verbatim}

In the case there are several parameter sets, add the name of the set for which you would like to the the parameters. In the ESPRESSO workflow there is only a parameter set. But in the case of KMOS, one can show the information for the parameter set: {\tt qc0\_parameters}, which is not the default one:

\begin{verbatim}
 edps -w kmos.kmos_wkf -p object qc0_parameters
\end{verbatim}

In general, to display simultaneously the information for all the parameter sets, type:

\begin{verbatim}
 edps -w espresso.espresso_wkf -ps
\end{verbatim}

\subsubsection{Use certain workflow and recipe parameter sets}
Each workflow can have different sets of parameters. There are two types of parameters: workflow parameters and recipe parameters. In the case of KMOS, there is a non-default parameter set named {\tt qc0\_parameters}. In order to use the workflow parameters from that set, type: 
\begin{verbatim}
 edps -w kmos.kmos_wkf -wps qc0_parameters
\end{verbatim}

In order to use the recipe parameters from that set, type: 
\begin{verbatim}
 edps -w kmos.kmos_wkf -rps qc0_parameters
\end{verbatim}


In order to use the both the workflow and recipe parameters from that set, type: 
\begin{verbatim}
 edps -w kmos.kmos_wkf -wps qc0_parameters -rps qc0_parameters
\end{verbatim}


\subsection{Configuration file for recipe and workflow parameters.}
\label{sec:edps_parameters_file}

Each workflow comes with a configuration file that contains recipe and workflow parameters. The recipe parameters are organised by task (indeed, several tasks can run the same recipe and might require different parameter values). If a recipe parameter is not listed, then the pipeline default is used. 
Each parameter file can contain more than one parameter set, therefore the users can define the set with the parameters that are most suited for their reduction. 


The configuration file is in {\tt yaml} format, therefore it must follow some conventions. Please use the default configuration file as starting point to create new ones. General rules that one has to keep in mind are:

\begin{itemize}
    \item Booleans has to be indicated as strings (e.g. "FALSE", "TRUE").
    \item Values has to be specified via column, not via equality. E.g.:\\ {\tt espresso.espdr\_sci\_red.cosmic\_detection\_sw: 2}
    \item Recipe parameters have to be specified by their full name, that follows the convention \\ {\tt <instrument>.<recipe>.<alias>}
\end{itemize}

Note that the direct specification of a parameter value in the {\tt edps} command overrides the values defined in the parameter file (see Section \ref{sec:edps_parameters}). 

\subsection{Use only calibration that delivers the desired product quality}
\label{sec:edps_quality}

By default the EDPS workflow tries to associate the closest in time suitable calibration to a dataset (e.g. a flat with the same filter). However, calibrations present in the input data directory can be too distant in time to ensure a certain quality of the results.
One can decide to associate only those calibrations that ensure the required quality of the final product. This is done by specifying the {\tt quality\_threshold} workflow parameter; EDPS will automatically select the most suited calibrations matching at least the specified quality threshold, considering that different instruments and different types of calibration might have different time validity range and conditions. 

To use only the calibrations that ensure a desired quality of the products, type on a terminal where the edps environment is active (single line command):

\begin{verbatim}
    edps -w espresso.espresso_wkf
         -i  <input_directories> 
         -wp quality_threshold n 
         -o <output_directory>
\end{verbatim}


where {\tt n} is a number with the following conventional meaning:
\begin{itemize}
    \item {\tt n<0}. Calibrations more restrictive than the calibration plan are selected.
    \item {\tt n=0}. Calibrations that follow the rules of the instrument calibration plans are selected.
    \item {\tt n=1}. The selected calibrations are sufficient to ensure good quality science results.
    \item {\tt n=2}. The selected calibrations "probably" produce results of still "acceptable" quality.
    \item {\tt n=3}. Significant risk of bad quality results or recipe failure.
\end{itemize}

By default {\tt n=inf}: if a suitable calibration is found, it is associated, regardless of the quality or the risk of having a pipeline crash.

In general, if all calibrations that should be provided by the standard calibration plan are present, it is not necessary to specify the quality level. This is the general case when downloading the data with CalSelector\footnote{\url{http://archive.eso.org/cms/application_support/calselectorInfo.html}}.

If the users want to be sure to reduce only datasets for which the calibrations from the calibration plan are provided, set {\tt n=0}. In this case, some datasets might be incomplete.


{\it Note:} During the association process in EDPS, the users can specify whether to give priority to {\tt RAW} calibrations (i.e. calibrations not yet processed by the pipeline) or to {\tt MASTER} calibrations (i.e. already reduced calibrations). This is explained in Section \ref{sec:edps_association_preference}.
